\documentclass[11pt]{article}

% ── Packages ─────────────────────────────────────────────────────────────
\usepackage[margin=1in]{geometry}
\usepackage{booktabs}
\usepackage{array}
\usepackage{longtable}
\usepackage{xcolor}
\usepackage{enumitem}
\usepackage{hyperref}
\usepackage[T1]{fontenc}

\hypersetup{colorlinks=true, linkcolor=blue!50!black, urlcolor=blue!50!black}

% A friendly "explain it simply" box used all over the document.
\newenvironment{simplebox}
  {\par\medskip\noindent\begin{quote}\color{blue!55!black}\small
   \textbf{\textsf{In simple words:}}\space\space}
  {\end{quote}\medskip}

% Shorthand for code/file names (underscores must be escaped by hand).
\newcommand{\code}[1]{\texttt{#1}}

\title{\textbf{AeroTwin -- Project Status Report}\\[4pt]
\large What we finished, what we fixed, and what is still left to do\\[2pt]
\normalsize (Written in plain language for someone brand-new to AI)}
\author{Team AeroTwin}
\date{September 12, 2026}

\begin{document}
\maketitle
\tableofcontents
\newpage

% ════════════════════════════════════════════════════════════════════════
\section{What Is This Project?}
% ════════════════════════════════════════════════════════════════════════

Big drones (the kind used by DRDO for long surveillance missions) use piston
engines -- the same basic idea as a car engine, but built to run for 18--24
hours straight in the sky. If that engine quietly starts to fail mid-mission,
a human watching screens might notice too late.

Our project, \textbf{AeroTwin}, is a computer program that:

\begin{enumerate}[nosep]
  \item \textbf{Fakes a drone engine} using physics equations (we cannot
        borrow a real engine, so we simulate one), sending out ``engine
        vitals'' like RPM, temperatures, oil pressure and vibration many
        times every second.
  \item \textbf{Watches those vitals} with AI programs that learned from
        examples what a healthy engine looks like and what different kinds
        of ``sickness'' (faults) look like.
  \item \textbf{Warns the human early} -- ideally about 30 minutes before a
        small problem becomes a dead engine -- and says \emph{why} it is
        worried, in plain language.
\end{enumerate}

\begin{simplebox}
Think of it like a doctor's smartwatch for a drone engine. The watch (our
software) constantly measures the heartbeat (RPM), temperature, and pulse.
It has studied thousands of days of health data, so when the pattern looks
``off'', it does not just beep -- it says: ``Your heart rate is rising
faster than normal for this activity, see a doctor soon.'' That last part
-- the \emph{explanation} -- is what makes our project more than a simple
alarm.
\end{simplebox}

\paragraph{A few words you will see everywhere}
\begin{description}[leftmargin=1.4em, style=nextline, font=\normalfont\bfseries]
  \item[Telemetry] The stream of measurements from the engine (RPM,
        temperatures, oil pressure, vibration\dots). Like the vitals monitor
        in a hospital.
  \item[Sensor] A small measuring device (a thermometer, a vibration meter).
  \item[Digital twin] A live, computer-game-style copy of the real engine
        running inside our software. The copy mirrors what the engine is
        doing right now.
  \item[Fault] Something going wrong (engine too hot, oil leaking, a sensor
        lying to us).
  \item[Fault injection] A button that makes our \emph{fake} engine sick on
        purpose, so we can test whether the AI notices. Essential for the
        demo.
  \item[Model / ML] A math program that learns patterns from examples.
        ``Training'' it means showing it lots of examples.
  \item[Anomaly] ``Something weird.'' Not a specific disease -- just a
        ``this does not look normal'' alarm.
  \item[RUL (Remaining Useful Life)] Our best guess of how much longer the
        engine (or part) can safely keep running.
\end{description}

% ════════════════════════════════════════════════════════════════════════
\section{How The Whole System Is Built}
% ════════════════════════════════════════════════════════════════════════

The system is a relay race. Data flows through four runners, left to right:

\begin{verbatim}
  [1] ENGINE SIMULATOR   ->  [2] AI BRAIN   ->  [3] SERVER  ->  [4] DASHBOARD
  (fake engine, sends      (checks the        (mailman:      (pretty screens
   vitals 10x/second)      vitals, spots      saves data     the judges see)
                           trouble,           + passes
                           guesses RUL)       messages)
\end{verbatim}

\begin{itemize}[nosep]
  \item \textbf{Runner 1 -- Simulator} (folder \code{simulator/}): physics
        equations for a 4-cylinder engine, mountain air and weather effects,
        and the fault-injection buttons.
  \item \textbf{Runner 2 -- AI services} (folders \code{ml/} and
        \code{fusion\_ml/}): cleans the data, compares measurements against
        what physics says they \emph{should} be, and runs the trained models
        (anomaly score, fault guess, wear level, RUL).
  \item \textbf{Runner 3 -- Servers} (folders \code{backend/} Python
        FastAPI + Node.js gateway): saves everything to a database and
        forwards messages.
  \item \textbf{Runner 4 -- Dashboard} (folder \code{frontend/}): a dark
        ``mission control'' website with gauges, charts, a rotating 3D
        engine, and replay of past missions.
\end{itemize}

\begin{simplebox}
There are actually \emph{two} sets of runners 1--3 in this codebase (an
older all-in-one Python version and a newer split-into-pieces version).
Both feed the same dashboard. For the report card we treat the AI training
pipeline (\code{ml/}) as the important one, because that is where our test
results come from.
\end{simplebox}

% ════════════════════════════════════════════════════════════════════════
\section{Report Card Words (How We Grade the AI)}
% ════════════════════════════════════════════════════════════════════════

Before reading the results, you need four grading words. These matter
because \textbf{accuracy alone can lie}.

\subsection*{Accuracy -- ``percent right'' (but it can fool you)}
If we test the AI 100 times and it answers right 86 times, accuracy is
86\,\%. Sounds great. But imagine a test where 90 questions are ``is the
engine healthy?'' and only 10 are about real faults. A lazy student who
answers ``healthy'' to \emph{everything} scores 90\,\% while catching
\emph{zero} faults. That lazy student is dangerous.

\subsection*{F1 score -- ``are you actually useful?'' (0 to 1, higher = better)}
F1 checks both: (a) when the AI shouts ``fault!'', how often is it right?
and (b) of all the real faults, how many did it catch? A lazy
``always healthy'' AI gets F1 near \textbf{0}. Our test 1 got
\textbf{0.1} -- that number told us the AI was (mostly) being lazy, even
with decent accuracy.

\subsection*{Confusion matrix -- ``show me exactly what you mixed up''}
A grid: rows = what was \emph{really} happening, columns = what the AI
\emph{guessed}. A perfect AI has all numbers on the diagonal. Numbers off
the diagonal show exactly which sickness it mistook for which -- like a
teacher showing which two answers you confused.

\subsection*{Anomaly score, degradation, RUL -- the three outputs}
\begin{itemize}[nosep]
  \item \textbf{Anomaly score} (0--100): how ``weird'' the engine looks
        right now. 0 = totally normal, 100 = very weird.
  \item \textbf{Degradation} (0--1): how \emph{worn out} the engine seems
        (0 = fresh, 1 = about to fail).
  \item \textbf{RUL}: estimated hours/days left before maintenance is due.
\end{itemize}

\begin{simplebox}
Accuracy is like counting how many baskets you scored, even if you never
played defense. F1 is the coach's honest grade: can you score \emph{and}
defend? That is why we report both, and why F1 matters more here -- in a
safety project, missing a real fault (a ``false negative'') is much worse
than one extra false alarm.
\end{simplebox}

% ════════════════════════════════════════════════════════════════════════
\section{Test 1: What Went Wrong}
% ════════════════════════════════════════════════════════════════════════

\subsection{The numbers we got}

\begin{center}
\begin{tabular}{lll}
\toprule
\textbf{Metric} & \textbf{Result} & \textbf{Verdict} \\
\midrule
Fault accuracy           & 86\,\%    & Looks nice but misleading (see below) \\
Fault F1 (macro)         & 0.1       & Very bad -- the AI was mostly lazy \\
Anomaly score quality    & poor      & One of our two worst metrics \\
RUL estimate quality     & poor      & The other worst metric \\
Top mix-up               & ALTERNATOR\_DEGRADATION & Most common wrong answer \\
\bottomrule
\end{tabular}
\end{center}

\subsection{Three sicknesses were never learned}
Three fault types --- \textbf{MISFIRE}, \textbf{LUBRICATION\_ISSUE}, and
\textbf{ABNORMAL\_VIBRATION} --- had so few example missions that the AI
basically never saw them during training. You cannot learn to recognize a
dog if you have only ever seen half a dog photo. They ``were not trained''.

\subsection{Why accuracy still said 86\,\%}
The test data was lopsided: far more ``healthy'' windows than ``fault''
windows. The AI learned the lazy trick --- ``just say healthy'' --- and
still collected a high score. The terrible F1 (0.1) exposed the trick.

\subsection{Why ALTERNATOR\_DEGRADATION was the top error}
The battery/alternator signals are \emph{small and quiet} compared to big,
loud signals like exhaust temperature. Their small numbers are easy to
blur into other fault patterns, so the AI kept blaming the alternator when
something else was wrong. (In our feature code, its ``noise budget'' is the
smallest of all channels, so tiny overlaps confuse the model.)

\begin{simplebox}
Test 1 summary: ``Our student got a B on a test where it copied the
answer `healthy' for most questions, never studied three chapters at all,
and its two real skills (spotting weirdness and predicting remaining life)
were the weakest subjects.''
\end{simplebox}

% ════════════════════════════════════════════════════════════════════════
\section{How We Fixed Test 1}
% ════════════════════════════════════════════════════════════════════════

Three changes, each aimed at one disease the test revealed:

\begin{enumerate}[nosep]
  \item \textbf{Balanced the data by fault type.} We evened out the
        examples so every sickness type got a fair number of practice
        questions. No more ``90 healthy vs 10 sick'' tests that reward
        laziness.
  \item \textbf{Added more missions.} More fake flights = more examples of
        each fault, so the three previously-untrained types finally had
        study material.
  \item \textbf{Added a confusion matrix.} Now we can literally see which
        sickness gets confused with which, instead of guessing from one
        overall grade.
\end{enumerate}

\subsection{Results after the fix}

\begin{center}
\begin{tabular}{lll}
\toprule
\textbf{Metric} & \textbf{Result} & \textbf{Reading} \\
\midrule
Anomaly value   & 0.096 & Very low false-alarm tendency (good) \\
Degradation     & 0.0   & Correct for a healthy engine \\
RUL estimate    & 500.3 h & Stable, no crazy numbers \\
Fault verdict   & none  & Correct -- the test run was healthy \\
\bottomrule
\end{tabular}
\end{center}

\begin{simplebox}
Important honesty note: this check ran on a \emph{healthy} engine. A clean
bill of health here proves we are \emph{not crying wolf}. It does not yet
prove we \emph{catch real sicknesses} -- that needs a test where a fault is
switched on and we watch the AI react. That is on the to-do list.
\end{simplebox}

% ════════════════════════════════════════════════════════════════════════
\section{Test 2: Making It Smarter}
% ════════════════════════════════════════════════════════════════════════

Test 2 attacked the remaining weak spots: anomaly detection, RUL, and fault
verdicts. Here is every change, explained:

\subsection{The seven upgrades}

\begin{description}[leftmargin=0em, style=nextline, font=\normalfont\bfseries]
  \item[1. Filter out ``transition'' windows.] When a drone changes speed
        or climbs, the engine briefly looks weird \emph{even when healthy}
        -- like you panting after sprinting up stairs. It is not sickness,
        it is effort. We now throw away those in-between moments so the AI
        does not study them as if they were faults.
  \item[2. FFT / frequency features.] ``FFT'' is a math trick that turns a
        wiggly line into its ``notes''. Instead of only knowing ``the
        engine shakes a bit'', we now know \emph{which rhythm} it shakes
        at. A loose part shakes at a different rhythm than a healthy
        engine -- like hearing a wobbly washing machine vs a balanced one.
  \item[3. Better window length and stride.] We changed how big each
        ``peek'' at the data is (the window) and how often we peek (the
        stride). Smaller, more frequent peeks catch problems sooner.
  \item[4. Temporal smoothing after prediction.] One noisy reading should
        not make the needle jump. We average the AI's answers over a short
        time so the display moves like a calm needle, not a frightened
        one.
  \item[5. Buffer masking.] At start-up the data buffers are half-empty
        and the numbers are garbage-ish. We now mask (ignore) that
        warm-up period instead of letting the AI learn from nonsense.
  \item[6. Simple extra stats (std, ptp, variance).] We added three cheap
        ``shape'' numbers per signal: standard deviation (how much it
        jitters), peak-to-peak (how far it swings), and variance (jitter
        squared -- especially for vibration). Cheap features, real value.
  \item[7. Rolling majority vote.] Instead of trusting a single answer, we
        take the AI's verdict over several consecutive windows and keep the
        most common one -- like asking the same question three times and
        going with the majority. One-off mistakes get outvoted.
\end{description}

\subsection{Results after Test 2}

\begin{center}
\begin{tabular}{lll}
\toprule
\textbf{Metric} & \textbf{Result} & \textbf{Reading} \\
\midrule
Anomaly score   & 67\,\%   & Big improvement on our weakest metric \\
Degradation     & 0.0      & Stable \\
RUL estimate    & 499.8 h  & Stable \\
Fault verdict   & none     & Stable \\
\bottomrule
\end{tabular}
\end{center}

\begin{simplebox}
Think of Test 2 as giving the student glasses (FFT rhythm-vision), a good
nights sleep (smoothing), a rule to skip gym class right after running
(transition filtering), and a friend to double-check answers (majority
vote). Same student, much better habits.
\end{simplebox}

% ════════════════════════════════════════════════════════════════════════
\section{The Big DONE / TODO Table}
% ════════════════════════════════════════════════════════════════════════

\begin{longtable}{p{4.6cm} p{2.0cm} p{7.6cm}}
\toprule
\textbf{Part of the system} & \textbf{Status} & \textbf{Notes (plain words)} \\
\midrule
\endhead
Fake engine (physics) & DONE & 4-cylinder engine with realistic temperatures, oil, fuel and vibration. \\
Sky \& weather effects & DONE & Air gets thinner with altitude; humidity/rain/dust tweak the engine. \\
Fault injection buttons & DONE & Can make the fake engine sick on purpose for demos. \\
\quad Unify fault name lists & TODO & The two code generations call faults by different names (e.g.\ one says \code{LEAN\_BURN\_RUNAWAY}, the other \code{INJECTOR\_DEGRADATION}). One shared list is needed so buttons and AI agree. \\
Telemetry streaming (10/s) & DONE & Engine vitals stream out live over the network. \\
FastAPI server (Python) & DONE & Saves missions, replays, reports; serves the dashboard. \\
Node gateway + TimescaleDB & DONE & Second-generation server with a proper time-series database. \\
Dashboard website & DONE & Dark mission-control UI: gauges, charts, alerts. \\
3D engine view & DONE & Rotating 3D engine colored by cylinder temperature. \\
Mission replay + reports & DONE & Rewatch a past flight and export a summary. \\
ML training pipeline & DONE & One command trains all models and prints a grade report. \\
Balanced training data & DONE & Test-1 fix: every fault type gets fair examples. \\
More missions & DONE & Test-1 fix: bigger dataset, previously-missing fault types now trained. \\
Confusion matrix & DONE & We can see exactly which faults get mixed up. \\
FFT / rhythm features & DONE & Test-2 fix: the AI can ``hear'' shake rhythms. \\
Window/stride tuning & DONE & Test-2 fix: smaller, more frequent peeks at the data. \\
Prediction smoothing & DONE & Test-2 fix: calm needle, no jumping. \\
Buffer masking & DONE & Test-2 fix: ignore warm-up garbage. \\
Extra stats (std/ptp/var) & DONE & Test-2 fix: cheap but useful signal-shape numbers. \\
Rolling majority vote & DONE & Test-2 fix: best-of-several answers. \\
\midrule
\textbf{Still to do} & & \\
\midrule
Commit trained model files & TODO & In this copy of the repo the model folder is \emph{empty} -- the trained models from tests 1--2 live in the newer working copy and must be saved into the project (so the demo never runs ``blind''). \\
Test RUL on a \emph{sick} engine & TODO & All our final numbers are from healthy runs. We must inject a fault and watch RUL actually count down. Until then, ``catches problems early'' is unproven. \\
Re-grade with per-fault F1 & TODO & After the fixes, re-run the full evaluation and report F1 for \emph{each} fault type, not just the average. This proves the three once-untrained types are now caught. \\
Confirm what ``anomaly 67\,\%'' means & TODO & Write down whether 67 is a score (0--100 weirdness) or a ``percent correct'' metric, so the numbers in the pitch match the code. \\
Explain RUL units change & TODO & The old code measured RUL in \emph{minutes-until-fault}; the new numbers are \emph{hours-until-maintenance} ($\approx$500\,h). Both are valid, but the presentation must say which one it is showing. \\
Demo rehearsal + fallbacks & TODO & Scripted 2-minute and 5-minute demos with screenshots as backup if live systems fail. \\
Real-engine calibration & TODO (future) & Our engine is physics equations, not a real DRDO engine. Real data would be needed before any real-world use. \\
\bottomrule
\end{longtable}

% ════════════════════════════════════════════════════════════════════════
\section{Honest Warnings (What Our Numbers Do and Do Not Prove)}
% ════════════════════════════════════════════════════════════════════════

\begin{enumerate}[nosep]
  \item \textbf{The engine is fake.} All data comes from our own physics
        simulator. That is perfect for learning and demos, but a real
        engine would be messier. We must always say ``demonstrated on
        synthetic telemetry''.
  \item \textbf{Our best numbers come from a healthy engine.} Degradation
        0.0, fault ``none'', stable RUL $\approx$ 500\,h -- all correct for
        a healthy run. They prove we do not cry wolf. They do not yet
        prove we catch real faults early. One faulted-run test would.
  \item \textbf{The two code generations must not be mixed up.} Fault
        names, message formats and ports differ between the old and new
        halves. For the demo, pick one path end-to-end (recommended: the
        new simulator $\rightarrow$ AI $\rightarrow$ gateway chain).
  \item \textbf{Model files must be treated like precious homework.} If
        the trained model files are not committed and the demo laptop
        loses them, the software silently falls back to a dumb rule-based
        mode. Always run a pre-demo check that models load.
\end{enumerate}

\begin{simplebox}
The golden rule for the pitch: \emph{under-promise, over-deliver.} Say
``early-warning capability, demonstrated on synthetic telemetry,
physics-informed prototype'' -- never ``certified'' or ``guaranteed to
prevent failure''. Judges trust honesty more than big claims.
\end{simplebox}

% ════════════════════════════════════════════════════════════════════════
\section{Next Steps (In Order of Importance)}
% ════════════════════════════════════════════════════════════════════════

\begin{enumerate}
  \item \textbf{Save the trained models into the repo} (model files +
        grade report + confusion matrix image). One hour of work, removes
        the biggest demo risk.
  \item \textbf{Run one faulted mission end-to-end.} Inject a fault, watch
        anomaly rise, RUL fall, alert fire. Screenshot everything -- this
        is the proof that our fixes actually catch sickness, not just stay
        quiet when healthy.
  \item \textbf{Re-run the full evaluation} and write down F1 for every
        fault type. Confirm MISFIRE, LUBRICATION and ABNORMAL\_VIBRATION
        are now caught.
  \item \textbf{Unify the fault-name lists} across simulator, AI and
        dashboard so the demo buttons match what the AI was trained on.
  \item \textbf{Write the two demo scripts} (2-minute backup, 5-minute
        full) and rehearse with the screenshots as fallback.
  \item \textbf{Update the pitch slides} with the Test-1 $\rightarrow$
        Test-2 improvement story: ``we found our AI was being lazy
        (accuracy 86\,\%, F1 0.1), fixed the data balance, added rhythm
        features and smoothing, and stabilized anomaly detection''.
        Judges love a team that finds its own bugs and fixes them.
\end{enumerate}

% ════════════════════════════════════════════════════════════════════════
\section{One-Page Summary (If You Read Nothing Else)}
% ════════════════════════════════════════════════════════════════════════

\begin{itemize}[nosep]
  \item We built a complete fake-drone-engine health monitor: fake engine
        $\rightarrow$ AI checks $\rightarrow$ servers $\rightarrow$
        mission-control dashboard, with replay and reports.
  \item \textbf{Test 1} exposed real problems: three fault types never
        trained, the AI was ``lazy'' (86\,\% accuracy but 0.1 F1), anomaly
        and RUL were the weakest skills, and the alternator was the
        favorite wrong answer.
  \item We fixed it by balancing the data, adding more missions, and
        adding a confusion matrix. Healthy-run results became clean:
        anomaly 0.096, degradation 0.0, RUL 500.3\,h, fault none.
  \item \textbf{Test 2} upgraded the AI's habits: skip transition moments,
        rhythm features (FFT), better windows, smoothing, buffer masking,
        extra stats, and majority voting. Anomaly performance improved to
        67\,\% with degradation 0.0 and RUL 499.8\,h.
  \item \textbf{Still open:} commit the trained models, prove RUL/anomaly
        on a \emph{sick} engine, re-grade per-fault F1, unify fault names,
        rehearse the demo.
  \item \textbf{Always say:} software demonstrator on synthetic telemetry
        -- an early-warning decision-support prototype, not a certified
        safety system.
\end{itemize}

\end{document}
